% =====================================================================================
% parzen-cdf project report -- self-contained version.
%
% STRUCTURE: one chapter per stage of the pipeline, in the order the data flows through
% them. Each chapter says what the stage does, why it is built that way, and what it
% replaces with respect to earlier versions of the project.
%
% This file is the skeleton: "TODO" comments say what belongs in each section.
% =====================================================================================

\documentclass[11pt]{article}

\usepackage[english]{babel}
\usepackage[T1]{fontenc}
\usepackage[a4paper,margin=1in]{geometry}
\usepackage{amsmath,amssymb,amsthm}
\usepackage{graphicx}
\usepackage{booktabs}
\usepackage{float}
\usepackage{xcolor}
\usepackage{hyperref}

\graphicspath{{../results/}}
\hypersetup{colorlinks=true, linkcolor=blue!50!black, urlcolor=blue!50!black}

\newcommand{\Fhat}{\widehat{F}}
\newcommand{\fhat}{\widehat{f}}
\newcommand{\shat}{\hat{\sigma}}
\newcommand{\hhat}{\hat{h}}
\newtheorem{theorem}{Theorem}
\newtheorem{proposition}[theorem]{Proposition}

\newcommand{\takeaway}[1]{\medskip\noindent\fcolorbox{black!30}{black!4}{%
  \parbox{0.97\linewidth}{\textbf{In short.} #1}}\medskip}

\title{Nonparametric estimation of a distribution function\\
and of the density derived from it\\
\large Parzen window and neural network: method, measurements, and limits}
\author{\texttt{parzen-cdf} project}
\date{\today}

\begin{document}
\maketitle

\begin{abstract}
% TODO: given a finite sample and no assumption on the distribution behind it, we estimate
% the CDF with a network trained on Parzen-window labels and differentiate it to obtain the
% density. Announce the three non-obvious choices: the window comes from cross-validation
% rather than from a rule of thumb; the architecture is a valid CDF by construction; the
% quality is reported through quantities computable without knowing the truth. Close with
% the limitation we state openly (sharp edges).
\end{abstract}

\tableofcontents
\newpage

% =====================================================================================
\section{Introduction}
% =====================================================================================

\subsection{The problem}
% TODO: x_1..x_n i.i.d. from an unknown F; estimate F and from it f. Why go through the CDF
% rather than estimating f directly: it is what the assignment asks, but it also has a
% technical reason (the CDF is smoother, and integration averages noise away).

\subsection{What is delivered}
% TODO: two functions evaluable at any point, not tables; a value of h1 in the required
% form h_n = h1/sqrt(n); a diagnostic block computable without the truth; a command line.

\subsection{How to read this document}
% TODO: every quantitative claim states how it was obtained. Proofs are in appendix A, the
% scripts producing each number in appendix B. The "what we do not claim" sections are
% deliberate.

\subsection{Where this work comes from}
% TODO: short and honest history; it explains why the project looks the way it does. First
% pass (window from Silverman's rule, mis-scaled for the logistic kernel); revision that
% fixed the scale and introduced the h1/sqrt(n) schedule with the 1.5-sigma rule; this
% third pass, which reopens the rule, the architecture and the way results are delivered.
% Say plainly what replaces what.

% =====================================================================================
\section{Development data: generating honest samples}
% =====================================================================================

\subsection{Why the generator matters even though it is not part of the pipeline}
% TODO: the samples we calibrate and measure on are produced by us. A biased generator
% would contaminate every downstream conclusion without giving any signal.

\subsection{Ancestral sampling from a mixture}
% TODO: a mixture defined as a two-stage process; the one-line proof that the method is
% exact; cost O(n).

\subsection{The Ziggurat algorithm}
% TODO: introduce it properly but briefly. The problem (uniforms to normals), the
% alternatives (inversion, Box-Muller, polar), the idea of equal-area strips, why it is
% exact rather than approximate. Make clear it is NumPy's code, not ours, and that it sits
% at a different level from ancestral sampling: the two compose, they are not alternatives.

\subsection{Checking the generator}
% TODO: the tests performed and their outcome. Stress the right test: uniformity of the
% p-values, not "how many rejections did I get". The check on consecutive seeds.

\subsection{A note on circularity}
% TODO: the argument "the generator never touches the CDF, so learning F is not cheating"
% does not hold, and should be replaced by the correct one: the estimator sees only the
% samples.

% =====================================================================================
\section{The Parzen window estimator}
% =====================================================================================

\subsection{Definition and intuition}
% TODO: a kernel centred on each observation; the average; the role of h.

\subsection{The logistic kernel and a closed-form CDF}
% TODO: why the logistic one: its antiderivative is the sigmoid, so the estimated CDF is
% exact and smooth, and the density is exactly its derivative. Comparison with the other
% kernels in the registry.

\subsection{The window must shrink as the sample grows}
% TODO: the consistency conditions (h to 0, n*h to infinity) and the h_n = h1/sqrt(n)
% schedule.

\subsection{How an estimate is scored}
% TODO: KS on the CDF, ISE on the density, and why the second is the one that matters
% here. The oracle as a normaliser, not as a selector. Relative efficiency.

% =====================================================================================
\section{The window: the central problem}
% =====================================================================================

\subsection{The rule of thumb, and why it could not work}
% TODO: the h1 = 1.5*sigma rule of the previous pass; the three methodological defects
% (self-referential benchmark, wrong anchoring quantity, wrong validation metric).

\subsection{An external test bench}
% TODO: the Marron and Wand densities, what they are and why they are used; our three
% additions.

\subsection{The selectors compared}
% TODO: list and form of each; the scale correction for the logistic kernel; the
% leave-one-density-out calibration for those with a free constant.

\subsection{The central result: the constant is not a constant}
% TODO: the spread of the optimal constant; the fact that recalibrating it makes things
% worse.

\subsection{Least-squares cross-validation, and its measured instability}
% TODO: the ranking; the instability measured, and why it does not transfer to the error.

\subsection{What the square-root-of-n schedule costs}
% TODO: the empirical exponent against the theoretical one; the measured cost of carrying
% h1 across budgets; why at a fixed budget it costs nothing.

\subsection{What is delivered}
% TODO: h1 estimated inside the required form; the clarification that h1 was never a
% universal constant.

% =====================================================================================
\section{The network: enforcing the shape of a distribution function}
% =====================================================================================

\subsection{What the output must guarantee}
% TODO: the four requirements (range, monotonicity, limits, differentiability) and why the
% fourth is needed here (the density is the derivative).

\subsection{The first attempt and its limits}
% TODO: MLP with a final sigmoid, monotonicity repaired downstream by a running maximum and
% a rescaling. What it actually guarantees and what it does not.

\subsection{The saturation theorem}
% TODO: statement and consequence: the tails are unreachable with finite parameters.
% Pointer to the proof. Supporting measurements.

\subsection{The chosen architecture}
% TODO: convex mixture of logistic CDFs; its relation to the Parzen estimator (same form,
% with learned weights, centres and windows); the reference to sigmoidal flows.

\subsection{The guarantees, proved}
% TODO: statement; the point that the hypotheses hold on the whole parameter space, not on
% a feasible region to be policed. Numerical verification.

\subsection{No loss of expressiveness}
% TODO: with J = n the family contains the Parzen estimator exactly; density of the family
% as J grows.

\subsection{Standardisation and equivariance}
% TODO: the measured problem (translate the data and the estimate collapses), the theorem,
% the fix.

\subsection{How many components}
% TODO: the calibration; the asymmetry of the risk; the value chosen.

% =====================================================================================
\section{Training}
% =====================================================================================

\subsection{The labels: the Parzen Neural Network recipe}
% TODO: leave-one-out labels at the sample points only; the formula and why the point's own
% contribution is subtracted.

\subsection{Why a deliberately narrow window}
% TODO: noisy but unbiased labels, averaged away by limited capacity.

\subsection{What is no longer needed}
% TODO: monotonicity penalty, rectification, clamping: why they were there and why with
% this architecture they have no object. The clamp broke the identity between density and
% derivative.

\subsection{Reproducibility}
% TODO: the defect found (the seed did not control initialisation) and how it was closed;
% the stronger property of the new architecture (no residual randomness at all).

% =====================================================================================
\section{From the distribution function to the density}
% =====================================================================================

\subsection{The closed-form derivative}
% TODO: the formula; the fact that the identity f = dF/dx is algebraic, not numerical.

\subsection{Why finite differences on a grid were a problem}
% TODO: dependence on resolution; the inconsistency introduced by the clamp.

% =====================================================================================
\section{Delivering an estimate without knowing the truth}
% =====================================================================================

\subsection{The evaluation domain}
% TODO: why a domain is needed; why it cannot come from the true distribution; the rule
% chosen and the criterion behind it; the principle (dispersion for the extent, window for
% the resolution).

\subsection{Diagnostics computable from the samples alone}
% TODO: which indicators inform and which mislead; the case of the comparison against the
% empirical CDF, which is the most natural trap.

\subsection{The interface}
% TODO: the single entry point and the command line; what it produces.

% =====================================================================================
\section{Results}
% =====================================================================================

\subsection{Against the previous version of the pipeline}
% TODO: the end-to-end table; the comment on the worst case; how to read the mass column.

\subsection{Against the Parzen estimator}
% TODO: when the network beats its teacher and when it does not.

\subsection{What the figures show}
% TODO: one or two figures with commentary.

% =====================================================================================
\section{Scope and limits}
% =====================================================================================

\subsection{Distributions from other families}
% TODO: the resilience study; the two suspicions refuted and the one confirmed.

\subsection{The structural limit}
% TODO: sharp edges; the check that raising capacity is not a free remedy; the decision to
% state the limit.

\subsection{What we do not claim}
% TODO: explicit list.

% =====================================================================================
\section{Conclusions}
% =====================================================================================
% TODO: what was built, the three choices that characterise it, and what would remain to do.

% =====================================================================================
\appendix
\section{Proofs}
% =====================================================================================
% TODO: T1 saturation; T2 the mixture is a CDF; T3 expressiveness; T5 the clamp adds mass;
% T6 equivariance.

% =====================================================================================
\section{Reproducibility}
% =====================================================================================
% TODO: table study -> script -> output file; how to re-run; environment.

% =====================================================================================
\section{Decision log}
% =====================================================================================
% TODO: compact table decision -> reason -> where it is measured.

\end{document}
